\subsection{Material transferred from the short note, September 2026}
\label{sec:257-transferred}

The September 2026 revision of the short note kept the sufficient conditions,
the exact reduction at $1/2$ and the exact descriptions of an arbitrary
rational-valued support, and moved the remaining checked statements here so
that no declaration left the corpus.  Every entry below names its source
declaration.  Nothing in this subsection is new mathematics.

\subsubsection*{Composite dilation and the foreign-divisor defect}

\begin{thm}[257tr:1 --- exact dilation identity]
For $A\subseteq\N$ and $a,x\in\Npos$ with $a\in A$, define the foreign-divisor
defect $\delta_A(a,x)=\#\{d\mid ax: d\in A,\ d\nmid x,\ d\ne a\}$.  Then the
divisor incidence satisfies
$\mathrm{sc}_A(ax)=\mathrm{sc}_A(x)+\mathbf 1_{a\nmid x}+\delta_A(a,x)$.
The hard point is that a composite multiplier can create support divisors
which are neither the multiplier nor old divisors of $x$; $\delta_A$ records
them.  If every member of $A$ is prime then $\delta_A(a,x)=0$ and the familiar
prime-support dilation formula is recovered.  For an
\texttt{OrthogonalPetalBouquet} witness $h_B$ and every $i$,
\[
\delta_A(\mathrm{ray}_{h_B}(i),x)\le
|h_B.\mathrm{exceptional}|+\mathrm{sc}_{\mathrm{range}(h_B.\mathrm{petal})}(x),
\]
because every non-exceptional foreign divisor injects into a unique petal
divisor of $x$.  The correction is already visible at $A=\{2,6\}$: multiplying
$x=1$ by $6$ creates the additional support divisor $2$, so
$\delta_A(6,1)=1$.
\par\noindent\ev{Lean} \scale{uniform} \coord{support-incidence}
\lean{supportCoeff\_mul\_eq\_add\_defect}{Erdos249257/CompositeDilationDefect.lean:30}
\lean{compositeDilationDefect\_eq\_zero\_of\_prime\_support}{Erdos249257/CompositeDilationDefect.lean:103}
\lean{supportCoeff\_mul\_prime\_support}{Erdos249257/CompositeDilationDefect.lean:119}
\lean{mem\_compositeDilationDefect\_bouquet\_imp\_exceptional\_or\_petal\_dvd}{Erdos249257/CompositeDilationDefect.lean:133}
\lean{compositeDilationDefect\_le\_exceptional\_add\_petalCoeff}{Erdos249257/CompositeDilationDefect.lean:151}
\lean{compositeDilationDefect\_two\_six\_fixture}{Erdos249257/CompositeDilationDefect.lean:218}
\end{thm}

\begin{obs}[257tr:2 --- the bouquet reduced-modulus identity]
For a bouquet $A=E\cup\{c_ip_i\}$ with $c_i\mid Q$, $e\mid Q$ for $e\in E$,
the $p_i>1$ pairwise coprime and coprime to $Q$, and $M_F=Q\prod_{j\in F}p_j$
for a finite index set $F$, one has $c_ip_i/\gcd(c_ip_i,M_F)=p_i$ for
$i\notin F$: a new ray retains its entire petal as the modulus governing its
hits along a progression with step $M_F$.  The alternating-core example
$Q=6$, $E=\varnothing$, $p_i=r_i^2$ over the primes $r_i$ after $3$, with
$c_i$ alternating between $2$ and $3$, is infinite.  The unconditional
irrationality of these supports follows from reciprocal summability alone and
needs no forced-slot selector, so this identity is structural information and
not a step in that proof.
\par\noindent\ev{Lean} \scale{uniform} \coord{bouquet}
\lean{ray\_div\_gcd\_frameStep}{Erdos249257/SupportSunflowerDichotomy.lean:192}
\lean{alternatingPrimeSquareSupport\_infinite}{Erdos249257/SupportSunflowerDichotomy.lean:337}
\end{obs}

\subsubsection*{Restricted subsum sets: the full construction}

For $J\subseteq\N$ let $\mathcal A_J$ be the subsum set generated by the
weights $w_{k+1}=(2^{k+1}-1)^{-1}$ over $k\in J$.  The short note keeps the
volume dichotomy and the injectivity of the selector; the construction behind
them is recorded here.

\begin{thm}[257tr:3 --- the restricted geometry, in full]
The restricted tail $T_J(n)=\sum_{k\in J}w_{n+k+1}$ is summable and satisfies
$T_J(n)<w_n$ for every $J$ and every $n\ge1$, because deleting weights only
shrinks a tail that already sits below $w_n$ at the full support.  The allowed
digit strings form a closed set; the digit map is injective on them; the
range is the restricted subsum set and agrees with the image description; the
restricted set is compact, closed, contained in $\mathcal A$, and therefore
nowhere dense; and for infinite $J$ the digit space is preperfect, so the
restricted set is perfect.  Changing one coordinate changes the value by
precisely its signed weight, so for $k\notin J$ adjoining $k$ splits the new
set into $\mathcal A_J$ and its disjoint translate by $w_{k+1}$, which doubles
the volume.  At $J=\N$ the restricted set is $\mathcal A$ itself.  The
division-free identity multiplies the face volume by $2^{|F|}$ for finite
$F$, monotonicity in $J$ traps a set with infinitely many forbidden
coordinates below finite-codimension faces of arbitrarily small dyadic
measure, and the two cases assemble into the dichotomy quoted in the short
note.
\par\noindent\ev{Lean} \scale{uniform} \coord{achievement-set}
\lean{selectedMersenneTail}{ErdosProblems/Erdos257/MersenneSubseriesRigidity.lean:20}
\lean{summable\_selectedMersenneTail}{ErdosProblems/Erdos257/MersenneSubseriesRigidity.lean:24}
\lean{SupportedMersenneDigits}{ErdosProblems/Erdos257/MersenneSubseriesRigidity.lean:45}
\lean{isClosed\_supportedMersenneDigitSet}{ErdosProblems/Erdos257/MersenneSubseriesRigidity.lean:66}
\lean{supportedMersenneAchievementSet\_eq\_image}{ErdosProblems/Erdos257/MersenneSubseriesRigidity.lean:79}
\lean{isClosed\_supportedMersenneAchievementSet}{ErdosProblems/Erdos257/MersenneSubseriesRigidity.lean:96}
\lean{supportedMersenneAchievementSet\_subset}{ErdosProblems/Erdos257/MersenneSubseriesRigidity.lean:103}
\lean{preperfect\_supportedMersenneAchievementSet}{ErdosProblems/Erdos257/MersenneSubseriesRigidity.lean:150}
\lean{summable\_mersenneDigitTerm}{ErdosProblems/Erdos257/MersenneSubseriesRigidity.lean:175}
\lean{positiveMersenneDigitValue\_update}{ErdosProblems/Erdos257/MersenneSubseriesRigidity.lean:182}
\lean{supportedMersenneAchievementSet\_univ}{ErdosProblems/Erdos257/MersenneSubseriesRigidity.lean:300}
\lean{volume\_supportedMersenneAchievementSet\_finset\_compl\_scaled}{ErdosProblems/Erdos257/MersenneSubseriesRigidity.lean:313}
\lean{supportedMersenneAchievementSet\_mono}{ErdosProblems/Erdos257/MersenneSubseriesRigidity.lean:358}
\end{thm}

\begin{obs}[257tr:4 --- an effective non-membership certificate]
Non-membership of a nonnegative target is visible at a single level, and for
a rational target strict failure is certified effectively: rational upper
bounds for the remaining tail converge to $T_n$, so once one lies below the
rational residual the fatal inequality is proved.  At $x=3/4$ the greedy
algorithm skips $w_1=1$ and leaves residual $3/4$, while the exact
zero-lookahead upper bound for the remaining tail is $2w_2=2/3<3/4$.
\par\noindent\ev{Lean} \scale{uniform} \coord{achievement-set}
\lean{three\_fourths\_not\_mem\_mersenneAchievementSet}{Erdos249257/GreedyAchievementSet.lean:1782}
\end{obs}

\subsubsection*{Squarefree divisor counts, and one method boundary}

\begin{obs}[257tr:5 --- the squarefree count]
The squarefree divisors of $n$ biject with the subsets of its prime factors,
so their number is $2^{\omega(n)}$ and the incidence of
$A_{\mathrm{sf}}=\{d\ge2:d\ \text{squarefree}\}$ is $2^{\omega(n)}-1$.  At
$n=12$ the squarefree divisors are $1,2,3,6$ and the incidence is $3$; at
$n=30$ the incidence is $7$.
\par\noindent\ev{Lean} \scale{uniform} \coord{squarefree}
\lean{squarefreeDivisors\_eq\_image}{ErdosProblems/Erdos257/SquarefreeSupportIncidence.lean:71}
\lean{card\_squarefreeDivisors}{ErdosProblems/Erdos257/SquarefreeSupportIncidence.lean:94}
\end{obs}

\begin{obs}[257tr:6 --- the prefix-lcm hypothesis has no full-support instance]
The denominator-gap criterion behind the factorial and power-of-two support
rows provably cannot reach the full support: the prefix lcm grows too slowly
for its hypothesis to hold.  So the settled families are not nested, and their
union does not exhaust the infinite supports.
\par\noindent\ev{Lean} \scale{uniform} \coord{certificate-kernel}
\lean{lcm\_gap\_hypothesis\_fails\_full\_support}{Erdos249257/CertificateKernel.lean:6460}
\end{obs}

\subsubsection*{Half-value routes retained from the short note}

\begin{thm}[257tr:7 --- positive skips characterise the half value]
The finite greedy core of $1/2$ turns into an exact local quotient row at
endpoint $2c-2$ whenever $c\ge4$ and the rational greedy remainder after
ranks through $c-1$ lies strictly between $0$ and $w_c$; the deficit below
$1/2$ is then below the next Mersenne weight, so a Boolean word fills the
entire upper-half capacity.  The residual cannot vanish at any finite rank,
because the odd denominator of every finite positive Mersenne sum rules out
exact finite exhaustion of $1/2$.  Consequently
\texttt{CofinalPositiveHalfGreedySkips} holds if and only if
$1/2\in\mathcal A$: proving cofinal positive skips already proves the
half-membership endpoint, and half-membership forces them.  Neither side is
established.
\par\noindent\ev{Lean} \scale{uniform} \coord{half-greedy}
\lean{greedyMersenneRemainderRat\_half\_pos}{Erdos249257/BooleanMobiusSkipRowCofinal.lean:32}
\lean{cofinalPositiveHalfGreedySkips\_iff\_half\_mem}{Erdos249257/BooleanMobiusSkipRowCofinal.lean:110}
\end{thm}

\begin{obs}[257tr:8 --- the two conditional producers, in their own coordinates]
A \texttt{HalfTerminalOnlyScaledVanishingSequence} places $1/2$ in the
achievement set by the vanishing normalised-carry estimate, and a cofinal
family of nonempty \texttt{CylinderStage}s converts into a weaker
terminal-only strip before the same closed-set step is applied.  Both
producers remain unconstructed; the short note records only their consequences.
\par\noindent\ev{Lean} \scale{uniform} \coord{half-greedy}
\lean{half\_mem\_mersenneAchievementSet\_of\_terminalScaledVanishing}{Erdos249257/TerminalOnlyScaledVanishing.lean:165}
\lean{cofinalTerminalOnlyStrip\_of\_cofinalCylinderStages}{Erdos249257/SuffixCylinderTerminalOnlyBridge.lean:290}
\end{obs}

\subsubsection*{The denominator-$21$ quotient-greedy frontier}

Let $r_N$ be the real greedy remainder of $1/21$ after rank $N$, $y_N=2^Nr_N$
and $c_{N+1}=2^{N+1}/(2^{N+1}-1)$.  The scaled recurrence is the nearly
doubling map sending $y_N$ to $2y_N-c_{N+1}$ when $c_{N+1}\le2y_N$ and to
$2y_N$ otherwise, so rank $N+1$ is skipped precisely when $2y_N<c_{N+1}$.

\begin{thm}[257tr:9 --- the rational lower-separatrix criterion]
For every nonnegative rational $q$, membership of $q$ in the achievement set
is equivalent to $2y_N(q)<2^{N+1}/(2^{N+1}-1)$ for cofinally many $N$; in
particular $1/21\in\mathcal A$ if and only if its scaled greedy orbit crosses
this moving lower separatrix beyond every cutoff.  The hard direction uses
the rational normal form: membership is equivalent to infinitely many omitted
exponents, and a skip is exactly a lower-separatrix crossing.  The equivalent
native form asks for one bounded cofinal return of $2^Nr_N$, with no
prescribed bound and no bounded return gaps.  Exact computation through rank
$200{,}000$ records $96$ zero-defect returns and $4{,}956$ returns to
$Q_N\le1$, with last observed ranks $193{,}690$ and $199{,}930$ and maximum
observed gap $492$; those finite data prove no cofinality, and their exact
boundary is recorded in the public result-atom
\texttt{erdos\_257:one\_over\_twenty\_one\_small\_lambert\_defect\_return\_evidence}.
\par\noindent\ev{Lean} \scale{uniform} \coord{greedy-orbit}
\lean{rat\_mem\_mersenneAchievementSet\_iff\_scaledLowerBranchCofinally}{Erdos249257/GreedyTrapDynamics.lean:152}
\lean{one\_div\_twentyOne\_mem\_iff\_scaledLowerBranchCofinally}{Erdos249257/GreedyTrapDynamics.lean:283}
\lean{one\_div\_twentyOne\_mem\_iff\_scaledRemainder\_cofinallyBounded}{Erdos249257/GreedyTrapDynamics.lean:275}
\lean{one\_div\_twenty\_one\_mem\_mersenneAchievementSet\_of\_defect\_le\_one\_cofinally}{Erdos249257/BooleanMobiusCarry.lean:2842}
\end{thm}

\begin{thm}[257tr:10 --- the fatal-branch quotient refinement]
Write $Q_N=\lfloor2^N/21\rfloor-P_N$ for the denominator-$21$ defect after the
six-periodic residue of $2^N\bmod21$ has been removed, where $P_N$ is the
integral numerator of the binary divisor-incidence prefix generated by the
greedy support; an exact identity relates $Q_N$, $2^Nr_N$ and a finite-prefix
divisor tail, so these are not independent models.  At even depth $2R$ let
$D_R$ be the Boolean support obtained by descending integer-greedy selection
on the exact scaled quotient weights and $s_R$ its terminal scalar remainder,
and let $\mathcal F_{21}$ be \texttt{TwentyOneFatalAlignedBranch}.  Then
$1/21\in\mathcal A$ if and only if $\mathcal F_{21}$ fails; an unbounded
sequence of ranks with $s_R\le2^R$ gives $1/21\in\mathcal A$; and on
$\mathcal F_{21}$ the quotient orbit is eventually strictly supercapacity and
follows the exact affine recurrence
$s_{R+1}=4s_R+\tau_R-\pi_R-(2^{R+1}+1)$ with $s_R>2^R$ and
$D_{R+1}=D_R\cup\{R+1\}$, where $\tau_R$ is the periodic target pulse and
$\pi_R$ the divisor pulse generated by $D_R$.  Every closed Boolean quotient
row is exactly the canonical quotient-greedy row, and an aligned crossing
from saturation into strict supercapacity forces a missing canonical ancestor
at one-third or two-thirds scale together with a real skipped exponent of
square-root-bounded defect.  These conclusions remove alternative late
Boolean branches; they do not exclude the full fatal, cofinite, aligned
branch.  Contradicting the recurrence would be sufficient for membership, and
the recurrence alone is not the exact membership endpoint, because synthetic
pulses can sustain it without being the actual divisor pulse.
\par\noindent\ev{Lean} \scale{uniform} \coord{quotient-greedy}
\lean{twentyOneEvenQuotientGreedySupport}{Erdos249257/TwentyOneQuotientGreedy.lean:32}
\lean{twentyOneEvenQuotientGreedyRemainder}{Erdos249257/TwentyOneQuotientGreedy.lean:39}
\lean{twentyOneClosedRow\_forces\_quotientGreedy}{Erdos249257/TwentyOneQuotientGreedy.lean:231}
\lean{twentyOneAlignedSaturatedCrossing\_forces\_canonical\_ancestor\_hole}{Erdos249257/TwentyOneQuotientGreedy.lean:5179}
\lean{twentyOneAlignedSaturatedCrossing\_forces\_scaled\_greedy\_skip}{Erdos249257/TwentyOneQuotientGreedy.lean:5223}
\lean{twentyOneGreedyDefect\_add\_mod\_div\_eq\_scaled\_remainder\_add\_tail}{Erdos249257/BooleanMobiusCarry.lean:2117}
\end{thm}

\begin{obs}[257tr:11 --- pointwise pulse bounds are too weak]
The exact six-step recurrence is
$Q_{N+6}=64Q_N+3(2^N\bmod21)-L_N$ with $L_N$ the actual weighted repair load.
The formerly plausible translation-invariant six-step contraction first fails
at $N=73$; what survives is the slope-aware repair-load equivalence.  So a
finite-memory invariant on the true orbit, if one exists, has to use the
correlated divisor pulses and not a pointwise bound.
\par\noindent\ev{Lean} \scale{uniform} \coord{quotient-greedy}
\lean{twentyOneGreedyDefect\_add\_six}{Erdos249257/BooleanMobiusCarry.lean:1892}
\lean{twentyOneGreedyDefect\_add\_six\_le\_div\_six\_iff\_repairLoad}{Erdos249257/BooleanMobiusCarry.lean:1927}
\end{obs}

\begin{obs}[257tr:12 --- the final-skip Diophantine signatures]
If $1/21\notin\mathcal A$, let $M$ be the last skipped exponent, $S_M$ its
finite skipped prefix and $a_M=1/21+\sum_{d\in S_M}(2^d-1)^{-1}$.  With
$E=\sum_{n\ge1}(2^n-1)^{-1}$ the fatal interval $0<a_M-E<\mathrm{gap}_M$ is an
exact one-sided approximation, and the reduced skipped-prefix denominator has
binary order equal to the lcm of the skipped exponents and hence at least
$M$.  Parity-sensitive gcd restrictions and adjacent-numerator product
inequalities are also checked.  These are necessary signatures; the missing
input is an upper-height or irrationality-measure estimate for $E$.
\par\noindent\ev{Lean} \scale{uniform} \coord{quotient-greedy}
\lean{lastTwentyOneSkip\_erdosBorwein\_fatalInterval}{Erdos249257/TwentyOneQuotientGreedy.lean:3735}
\lean{lastTwentyOneSkip\_oddDoublingOrder\_eq\_lcm}{Erdos249257/TwentyOneQuotientGreedy.lean:3569}
\lean{lastTwentyOneSkip\_le\_oddDoublingOrder}{Erdos249257/TwentyOneQuotientGreedy.lean:3583}
\end{obs}

\begin{obs}[257tr:13 --- the primitive cone at $1/21$]
Write $\mathcal P_{23}(n)=\{(p,q)\in\Npos^2:\gcd(p,q)=1,\ 2p+3q=n\}$.  Every
$n\ge11$ has a member, $\mathcal P_{23}(10)$ is empty, multiplicity begins
immediately at $11$, and every $10k$ with $k\ge2$ has at least two distinct
primitive representations.  These are exact Diophantine facts.  The recurring
collisions show why primitive-cone coverage is not already a Boolean
expansion: the missing step is a proof that the relevant integer
multiplicities can be converted into zero-one reciprocal-Mersenne digits
without changing the value.
\par\noindent\ev{Lean} \scale{uniform} \coord{quotient-greedy}
\lean{no\_primitive23\_solution\_ten}{Erdos249257/Primitive23Multiplicity.lean:26}
\lean{exists\_two\_primitive23\_solutions\_eleven}{Erdos249257/Primitive23Multiplicity.lean:38}
\lean{exists\_primitive23\_solution\_of\_eleven\_le}{Erdos249257/Primitive23Multiplicity.lean:52}
\lean{exists\_two\_primitive23\_solutions\_mul\_ten}{Erdos249257/Primitive23Multiplicity.lean:86}
\end{obs}

\subsubsection*{Remaining questions carried over}

The short note keeps three endpoint questions: the universal problem on
reciprocal-divergent supports, membership of $1/2$, and membership of
$1/21$.  Four further questions belong with the material above.  Can the
actual denominator-$21$ greedy orbit satisfy $\mathcal F_{21}$ indefinitely,
or must an arbitrarily deep closed return $s_R\le2^R$ occur?  Is there a
finite-memory, $2$-adic or discrepancy invariant, built from a bounded window
of $R\bmod6$, residues of $s_R$, endpoint divisor counts and the finite set
of eventual skips, that forces descent or a forbidden state, or can one prove
that no bounded-memory invariant distinguishes the true orbit from synthetic
permanent-supercapacity controls?  Can the complete final-skip signatures of
257tr:12 be used to prove $|E-a_M|\ge\mathrm{gap}_M$?  And for $L\ge1$,
classify the reduced rationals $p/q\in(0,1)$ with $q$ odd,
$\operatorname{ord}_q(2)\le L$ and no finite Mersenne representation, for
which non-membership reduces to an eventually deterministic finite-memory or
affine quotient recurrence; in particular classify the case $L=6$ by
target-pulse alphabet, finite-support obstruction, weakest membership
criterion, surviving branch complexity and available Diophantine estimate,
and decide whether $1/21$ is minimal or unique under explicit criteria.
